\clearpage
\begingroup
\fancyhead[L]{\small NACHLASS UND BRIEFWECHSEL}
\renewcommand{\headrulewidth}{0pt}
\par\medskip
\SourceLeaf{589}
\Needspace{8\baselineskip}
\pdfbookmark[2]{[4] Nicolai an Gauss, 22. April 1814}{pallas-VIII-letter4}
\begin{center}
[4.]\par\medskip
\textsc{Nicolai} an \textsc{Gauss}, Seeberg, 22. April 1814.
\end{center}
\par\noindent\hspace*{\parindent}\{.\ .\ .\ .\ .\quad Meinen letzten Brief vom 8. April mit den Störungsformeln für
Excentricität und Perihelium werden Sie bekommen haben; heute kann ich
Ihnen das Ende der Arbeit schicken, nemlich die Gleichungen für den 2\textsuperscript{ten}
Theil der Epoche.\quad .\ .\ .\ .\ .

\par\noindent\hspace*{\parindent}Bei dem analytischen Ausdrucke für \(\mathrm{d}L\) habe ich eine kleine Transformation
in Anwendung gebracht, die ich Ihnen doch zur gütigen Ansicht mittheilen will.

\par\noindent\hspace*{\parindent}Setzt man nemlich die Grösse \(1+\cos v\cos E=u\), so wird
\[
\begin{aligned}
\mathrm{d}\varphi&=-maa\sin v\cos E\mathord{.}\,X+maauY\\
\mathrm{d}\tilde\omega&=A-maa\operatorname{cotang}\varphi\mathord{.}\,(3-u)\mathord{.}\,X
  +\frac{maa}{e}\cos v\sin E\mathord{.}\,Y.
\end{aligned}
\]

\par\noindent\hspace*{\parindent}Wird nun letzteres, mit Ausschluss des \(A\), mit \(1-\dfrac{rr}{aa\cos\varphi}\), ersteres mit
\(-\left(\cos\varphi+\dfrac{r}{a\cos\varphi}\right)\sin E\) multiplicirt, so folgt unmittelbar
\[
\begin{aligned}
\mathrm{d}L={}&A+\frac{m}{e}\bigl\{(rr-aa\cos\varphi)(3-u)+ae(p+r)(2-u)\cos E\bigr\}X\\
 &{}-\frac{m\sin E}{e\cos\varphi}\bigl\{(rr-aa\cos\varphi)\cos v+ae(p+r)u\bigr\}Y.
\end{aligned}
\]

\par\noindent\hspace*{\parindent}Ich habe es nicht versucht, die eingeklammerten Factoren noch mehr zusammenzuziehen,
weil sie mir zur numerischen Berechnung, zumal mit Ihrer
Hülfstafel, bequem genug schienen. Diese Umwandlung und Einführung der
Grösse \(u\) gründet sich übrigens auf eine Eigenschaft der Relation zwischen
der wahren und excentrischen Anomalie, auf welche ich vor Kurzem zufällig
gerieth. Denkt man sich nemlich die Complemente von \(v\) und \(E\) zu \(90^0\)
(oder zu \(450^0\)) als die beiden Seiten eines sphärischen Dreiecks, welche den
Excentricitätswinkel \(\varphi\) einschliessen, so ist die dritte, dem \(\varphi\) gegenüberstehende,
Seite dieses Dreiecks immer constant und zwar ebenfalls \(=\varphi\). Auch die
beiden andern Winkel dieses Dreiecks haben sonderbare Eigenthümlichkeiten;
der der Seite \(\genfrac{}{}{0pt}{0}{90^0}{450^0}-v\) gegenüberstehende ist dieser Seite immer gleich, und
der andere, der Seite \(\genfrac{}{}{0pt}{0}{90^0}{450^0}-E\) entgegengesetzte, ist das Complement von dieser
Seite zu \(180^0\) oder \(=90^0+E\).\quad .\ .\ .\ .\ .

\SourceLeaf{590}
\par\noindent\hspace*{\parindent}Sollten wir bei den Factoren der störenden Kräfte nicht ganz genau
stimmen, so wird dies daher rühren, dass Sie wahrscheinlich die neuesten
mittlern Elemente zum Grunde gelegt haben, ich hingegen diejenigen, welche
die ganz erste Manipelrechnung gegeben hat.
\par\medskip
\Needspace{13\baselineskip}
\begin{center}
Zweiter Theil der Epoche.\par
Jährliche Änderung \(=-1\overset{\prime\prime}{,}31973\)
\end{center}
\begingroup
\small
\newcommand{\VIIIpallas}{\text{\AstroSymbols ⚴}}
\newcommand{\VIIIsaturn}{\text{\AstroSymbols ♄}}
\setlength{\arraycolsep}{2pt}
\renewcommand{\arraystretch}{1.25}
\[
\begin{array}{r@{\,}l@{\qquad}r@{\,}l}
0\overset{\prime\prime}{,}45&\sin(\VIIIpallas+180^0\ 35')&+0\overset{\prime\prime}{,}27&\sin(3\VIIIpallas-2\VIIIsaturn+65^0\ 10')\\
+0{,}29&\sin(2\VIIIpallas+181\ 47)&+0{,}30&\sin(3\VIIIsaturn+207\ 41)\\
+1{,}78&\sin(\VIIIsaturn+140\ 25)&+1{,}05&\sin(\VIIIpallas-3\VIIIsaturn+163\ 1)\\
+0{,}32&\sin(\VIIIpallas+\VIIIsaturn+181\ 12)&+0{,}45&\sin(2\VIIIpallas-3\VIIIsaturn+154\ 50)\\
+0{,}41&\sin(\VIIIpallas-\VIIIsaturn+178\ 5)&+0{,}42&\sin(3\VIIIpallas-3\VIIIsaturn+173\ 25)\\
+0{,}11&\sin(2\VIIIpallas-\VIIIsaturn+158\ 28)&+0{,}21&\sin(\VIIIpallas-4\VIIIsaturn+166\ 29)\\
+2{,}01&\sin(2\VIIIsaturn+221\ 24)&+0{,}15&\sin(2\VIIIpallas-4\VIIIsaturn+159\ 55)\\
+3{,}35&\sin(\VIIIpallas-2\VIIIsaturn+124\ 44)&+0{,}11&\sin(3\VIIIpallas-4\VIIIsaturn+156\ 39).\}\\
+2{,}04&\sin(2\VIIIpallas-2\VIIIsaturn+178\ 31)&&
\end{array}
\]
\endgroup
\clearpage
\endgroup
